\documentclass{csse4400}

% \teachermodetrue

\usepackage{float}
\usepackage{languages}

\title{Managed Authentication with Amazon Cognito}
\author{Cameron Badman}

\date{\week[practical]{10}}
\begin{document}

\maketitle

\aside{
  This practical is intended to be a short authentication clinic before project team meetings.
}

\section{This Week}
Our goal is to add managed authentication to TaskOverflow without making TaskOverflow responsible for storing passwords.
Specifically, we will:
\begin{itemize}
  \item create an Amazon Cognito User Pool using Terraform;
  \item sign in as a test user and retrieve a JSON Web Token (JWT);
  \item add token verification code to TaskOverflow; and
  \item protect one TaskOverflow endpoint with an authentication decorator.
\end{itemize}

\section{How This Works}

Up to this point, TaskOverflow has focused on API design, persistence, deployment, scaling, events, and observability.
One important concern is still missing: user identity.
Rather than building password storage and account recovery ourselves, we will use Amazon Cognito as a managed identity provider.

\begin{description}
  \item[Cognito] owns users, passwords, and sign-in.
  \item[TaskOverflow] receives signed tokens and checks whether they are valid.
  \item[Route handlers] focus on business logic after authentication has succeeded.
\end{description}

The request path is:
\begin{enumerate}
  \item A user signs in with Cognito.
  \item Cognito returns a signed JWT.
  \item The client sends that token to TaskOverflow in the \texttt{Authorization} header.
  \item TaskOverflow verifies the token signature and claims.
  \item If the token is valid, the request continues to the protected route.
\end{enumerate}

\info{
  This practical deliberately avoids hosted login pages, front-end integration, MFA, and groups.
  Those are useful production features, but they are outside a 45 minute practical.
}

\section{Creating a Cognito User Pool}

Create a new file called \texttt{auth.tf} in the Terraform directory for your TaskOverflow deployment.
Add the following resources.

\begin{code}[language=terraform,numbers=none]{auth.tf}
resource "aws_cognito_user_pool" "taskoverflow" {
  name = "taskoverflow-auth"

  username_attributes      = ["email"]
  auto_verified_attributes = ["email"]

  password_policy {
    minimum_length    = 8
    require_lowercase = true
    require_numbers   = true
    require_symbols   = false
    require_uppercase = true
  }
}

resource "aws_cognito_user_pool_client" "taskoverflow" {
  name         = "taskoverflow-cli"
  user_pool_id = aws_cognito_user_pool.taskoverflow.id

  generate_secret = false

  explicit_auth_flows = [
    "ALLOW_USER_PASSWORD_AUTH",
    "ALLOW_REFRESH_TOKEN_AUTH"
  ]
}

output "cognito_user_pool_id" {
  value       = aws_cognito_user_pool.taskoverflow.id
  description = "Cognito User Pool ID for TaskOverflow."
}

output "cognito_client_id" {
  value       = aws_cognito_user_pool_client.taskoverflow.id
  description = "Cognito App Client ID for TaskOverflow."
}
\end{code}

Apply the Terraform configuration.

\bash{terraform apply}

\teacher{
  The app client does not have a secret because we want students to authenticate from the command line.
  A public browser or mobile app also cannot safely keep a client secret.
}

\section{Creating a Test User}

Save the values Terraform printed.

\begin{code}[language=bash,numbers=none]{}
export COGNITO_USER_POOL_ID=$(terraform output -raw cognito_user_pool_id)
export COGNITO_CLIENT_ID=$(terraform output -raw cognito_client_id)
export AWS_REGION=us-east-1
\end{code}

Create a user for testing.
Replace the email address with your own.

\begin{code}[language=bash,numbers=none]{}
aws cognito-idp admin-create-user \
  --user-pool-id "$COGNITO_USER_POOL_ID" \
  --username "student@example.com" \
  --user-attributes Name=email,Value="student@example.com" Name=email_verified,Value=true \
  --message-action SUPPRESS

aws cognito-idp admin-set-user-password \
  --user-pool-id "$COGNITO_USER_POOL_ID" \
  --username "student@example.com" \
  --password "Taskoverflow1" \
  --permanent
\end{code}

\warning{
  Do not use this password pattern in a real system.
  It is only here to keep the practical short and repeatable.
}

Now sign in and store the ID token in a shell variable.

\begin{code}[language=bash,numbers=none]{}
export ID_TOKEN=$(aws cognito-idp initiate-auth \
  --auth-flow USER_PASSWORD_AUTH \
  --client-id "$COGNITO_CLIENT_ID" \
  --auth-parameters USERNAME="student@example.com",PASSWORD="Taskoverflow1" \
  --query "AuthenticationResult.IdToken" \
  --output text)
\end{code}

\noindent
The token is what TaskOverflow will receive in the \texttt{Authorization} header.
It is signed by Cognito and includes claims about the authenticated user.

\section{Verifying Tokens in TaskOverflow}

Install the JWT dependency.

\bash{pipx run poetry add "pyjwt[crypto]"}

\noindent
Create \texttt{todo/auth/cognito.py}.
This file verifies tokens issued by the Cognito User Pool.

\begin{code}[language=python,numbers=none]{todo/auth/cognito.py}
import os

import jwt
from jwt import PyJWKClient


class AuthError(Exception):
    pass


def _required_env(name):
    value = os.environ.get(name)
    if not value:
        raise AuthError(f"Missing environment variable: {name}")
    return value


def verify_cognito_token(token):
    region = _required_env("AWS_REGION")
    user_pool_id = _required_env("COGNITO_USER_POOL_ID")
    client_id = _required_env("COGNITO_CLIENT_ID")

    issuer = f"https://cognito-idp.{region}.amazonaws.com/{user_pool_id}"
    jwks_url = f"{issuer}/.well-known/jwks.json"

    try:
        signing_key = PyJWKClient(jwks_url).get_signing_key_from_jwt(token)
        claims = jwt.decode(
            token,
            signing_key.key,
            algorithms=["RS256"],
            audience=client_id,
            issuer=issuer,
        )
    except jwt.PyJWTError as err:
        raise AuthError("Invalid authentication token") from err

    if claims.get("token_use") != "id":
        raise AuthError("Expected an ID token")

    return claims
\end{code}

\section{Adding an Authentication Decorator}

Now create \texttt{todo/auth/decorators.py}.
This decorator reads the \texttt{Authorization} header before the route handler runs.

\begin{code}[language=python,numbers=none]{todo/auth/decorators.py}
from functools import wraps

from flask import g, jsonify, request

from .cognito import AuthError, verify_cognito_token


def _bearer_token():
    header = request.headers.get("Authorization", "")
    prefix = "Bearer "
    if not header.startswith(prefix):
        raise AuthError("Missing bearer token")
    return header[len(prefix):]


def require_auth(view):
    @wraps(view)
    def wrapped(*args, **kwargs):
        try:
            g.user = verify_cognito_token(_bearer_token())
        except AuthError as err:
            return jsonify({"error": str(err)}), 401
        return view(*args, **kwargs)

    return wrapped
\end{code}

Create an empty \texttt{todo/auth/\_\_init\_\_.py} file so Python treats the folder as a package.

\section{Protecting a Route}

Add a route that allows us to check which user is authenticated.
In \texttt{todo/views/routes.py}, import the decorator and Flask's \texttt{g} object.

\begin{code}[language=python,numbers=none]{todo/views/routes.py}
from flask import g
from todo.auth.decorators import require_auth
\end{code}

Then add the following route.

\begin{code}[language=python,numbers=none]{todo/views/routes.py}
@api.route("/auth/me", methods=["GET"])
@require_auth
def me():
    return jsonify({
        "sub": g.user.get("sub"),
        "email": g.user.get("email"),
        "username": g.user.get("cognito:username"),
    }), 200
\end{code}

Finally, protect one existing write endpoint.
For example, if your TaskOverflow route for creating todos looks like this:

\begin{code}[language=python,numbers=none]{todo/views/routes.py}
@api.route("/todos", methods=["POST"])
def create_todo():
    ...
\end{code}

\noindent
add the decorator.

\begin{code}[language=python,numbers=none]{todo/views/routes.py}
@api.route("/todos", methods=["POST"])
@require_auth
def create_todo():
    ...
\end{code}

\teacher{
  The exact route name may differ depending on the starter repository version.
  The important point is that students protect one route by adding a single authentication decorator.
}

\section{Testing the Request Path}

Start TaskOverflow with the Cognito environment variables.

\begin{code}[language=bash,numbers=none]{}
export AWS_REGION=us-east-1
export COGNITO_USER_POOL_ID=$(terraform output -raw cognito_user_pool_id)
export COGNITO_CLIENT_ID=$(terraform output -raw cognito_client_id)

pipx run poetry run flask --app todo run --port 6400 --debug
\end{code}

\noindent
If you are running TaskOverflow with Docker Compose, add the same values to the app service environment instead.

\begin{code}[numbers=none]{docker-compose.yml}
    environment:
      AWS_REGION: us-east-1
      COGNITO_USER_POOL_ID: ${COGNITO_USER_POOL_ID}
      COGNITO_CLIENT_ID: ${COGNITO_CLIENT_ID}
\end{code}

\subsection{Unauthenticated Request}

Call the new endpoint without a token.

\bash{curl -i http://localhost:6400/api/v1/auth/me}

\noindent
You should receive a \texttt{401 Unauthorized} response.

\subsection{Authenticated Request}

Call the endpoint with the token.

\begin{code}[language=bash,numbers=none]{}
curl -i \
  -H "Authorization: Bearer $ID_TOKEN" \
  http://localhost:6400/api/v1/auth/me
\end{code}

\noindent
This request should return information from the Cognito token claims.

\section{Reflection}

For your project, consider:
\begin{itemize}
  \item Which endpoints should be public?
  \item Which endpoints require an authenticated user?
  \item Which endpoints require a specific role or group?
  \item Where should authentication be enforced: API Gateway, application middleware, or both?
\end{itemize}

\warning{
  Remember to destroy AWS resources when you are finished.
}

\bash{terraform destroy}

\bibliographystyle{ieeetr}
\bibliography{books,ours}

\end{document}
